\secnumbersection{MARCO CONCEPTUAL}

Este capítulo presenta los fundamentos de la modelación de flujos origen-destino. Describe los modelos clásicos como antecedentes teóricos y expone la formulación de \textit{Deep Gravity}, distinguiendo la publicación original de la adaptación utilizada en Santiago. Luego introduce los métodos de atribución de predicciones, las fuentes geoespaciales y el contexto de movilidad. El procedimiento experimental de esta memoria se especifica en el Capítulo~\ref{sec:cap4}.

% ----------------------------------------------------------------
\subsection{Movilidad urbana y flujos origen-destino}
% ----------------------------------------------------------------

La movilidad urbana comprende el conjunto de desplazamientos que realizan las personas entre distintas ubicaciones de la ciudad en un período de tiempo determinado \cite{liu2024interdisciplinary}. Comprender estos desplazamientos es un insumo fundamental para la planificación del transporte, el diseño de infraestructura vial y la formulación de políticas públicas basadas en evidencia.

Una representación sintética y cuantitativa de la movilidad urbana es la matriz origen-destino (matriz OD). Formalmente, dado un conjunto de $n$ zonas geográficas $\mathcal{Z} = \{z_1, z_2, \dots, z_n\}$ que particionan el área de estudio, el flujo $T_{ij}$ representa el número de viajes que se realizan desde la zona de origen $z_i$ hasta la zona de destino $z_j$ en un intervalo de tiempo fijo:
\begin{equation}
    \mathbf{T} \in \mathbb{R}^{n \times n}, \quad T_{ij} \geq 0 \quad \forall\, i, j \in \{1, \dots, n\}
    \label{eq:matriz_od}
\end{equation}
La colección de todos los flujos posibles entre pares de zonas conforma la matriz OD. Cada fila $i$ representa los flujos salientes desde una zona de origen, y cada columna $j$ los flujos entrantes hacia una zona de destino.

Obtener la matriz OD mediante observación directa, a través de encuestas domiciliarias, aforos de tráfico o registros de tarjetas de pago, es costoso y poco frecuente. Esta limitación ha motivado el desarrollo de modelos generativos de flujos, capaces de estimar $\mathbf{T}$ a partir de características observables del territorio, como la distribución de la población o la presencia de equipamientos urbanos, sin depender exclusivamente de mediciones directas \cite{simini2021deep}.

El problema de generación de flujos se formula de la siguiente manera: dado un conjunto de zonas $\mathcal{Z}$, sus características geoespaciales y los flujos totales salientes por zona $O_i = \sum_j T_{ij}$, estimar los flujos por par origen-destino $\hat{T}_{ij}$ para todos los pares $(i, j)$.

% ----------------------------------------------------------------
\subsection{Modelos clásicos de generación de flujos}
% ----------------------------------------------------------------

La estimación de flujos origen-destino cuenta con una tradición extensa en la geografía cuantitativa. Los modelos que se describen a continuación constituyen antecedentes teóricos de esta memoria. Los modelos de referencia utilizados en la evaluación de Santiago se especifican en el Capítulo~\ref{sec:cap4}.

\subsubsection{Modelo gravitatorio}

El modelo gravitatorio de movilidad establece una analogía con la ley de gravitación universal de Newton. La aplicación de este principio a los fenómenos sociales cuenta con antecedentes desde el siglo XIX. Posteriormente, Zipf \cite{zipf1946p} lo formalizó en 1946 para la predicción de flujos. Su formulación postula que el volumen de desplazamientos entre dos zonas es directamente proporcional al producto de sus masas (habitualmente la población) e inversamente proporcional a la distancia que las separa:
\begin{equation}
    T_{ij} = k \cdot \frac{P_i^{\,\alpha} \cdot P_j^{\,\beta}}{d_{ij}^{\,\gamma}}
    \label{eq:gravedad}
\end{equation}
donde $P_i$ y $P_j$ son las poblaciones de las zonas de origen y destino, $d_{ij}$ es la distancia entre sus centroides, $k$ es una constante de proporcionalidad, y $\alpha$, $\beta$, $\gamma$ son parámetros que deben calibrarse con datos observados. Zipf verificó empíricamente esta relación mediante datos de tráfico de buses, ferrocarriles y llamadas telefónicas entre ciudades de Estados Unidos. Los coeficientes obtenidos se aproximan a la unidad en todos los casos.

La formulación de la ecuación \eqref{eq:gravedad} recibe el nombre de \textit{modelo de gravedad no restringido}. En la práctica, para garantizar que los flujos generados sean consistentes con los totales salientes observados, se utiliza la versión \textit{restringida en origen} (\textit{singly-constrained}), que normaliza los flujos de cada origen sobre todos sus destinos posibles:
\begin{equation}
    \hat{T}_{ij} = O_i \cdot \frac{P_j^{\,\beta}\, f(d_{ij})}{\displaystyle\sum_{k} P_k^{\,\beta}\, f(d_{ik})}
    \label{eq:gravedad_restringida}
\end{equation}
donde $O_i = \sum_j T_{ij}$ es el flujo total saliente de la zona $i$ y $f(d_{ij})$ es una función de disuasión por distancia, que puede adoptar una forma exponencial, $f(d_{ij})=e^{-\lambda d_{ij}}$, con $\lambda>0$, o una ley de potencia, $f(d_{ij})=d_{ij}^{-\gamma}$, con $\gamma>0$ y $d_{ij}>0$\footnote{Existe también la versión \textit{doblemente restringida} (\textit{doubly-constrained}), que impone conservación tanto de los flujos salientes $O_i$ como de los flujos entrantes $D_j = \sum_i T_{ij}$. \textit{Deep Gravity} utiliza exclusivamente la versión restringida en origen.}.

Las limitaciones de un modelo gravitatorio dependen de su especificación \cite{liu2024interdisciplinary}. La versión presentada relaciona población, distancia y flujo mediante parámetros calibrados con observaciones. Al no incluir otros atributos, no distingue de forma directa usos del suelo, equipamientos ni infraestructura de transporte. Estas restricciones corresponden a esa formulación, no a toda extensión posible de la familia gravitatoria. Su adecuación y la transferencia de sus parámetros requieren evaluación empírica.

\subsubsection{Modelo de oportunidades intervinientes}

Stouffer \cite{stouffer1940intervening} propuso en 1940 una perspectiva alternativa al modelo gravitatorio, argumentando que la distancia no es en sí misma el factor determinante de la movilidad, sino que actúa como un contenedor espacial de oportunidades. Su hipótesis central es que el número de personas que se desplazan a una zona de destino es directamente proporcional al número de oportunidades en ese destino e inversamente proporcional al número de oportunidades intervinientes entre el origen y el destino.

Formalmente, sea $y$ el número de personas que se desplazan desde un origen hasta una banda circular de anchura $\Delta s$, y sea $x$ el número acumulado de oportunidades entre el origen y la distancia $s$. Stouffer postula:
\begin{equation}
    \Delta y = a \cdot \frac{\Delta x}{x}
    \label{eq:stouffer}
\end{equation}
Integrando esta expresión diferencial, el número total de personas que se detienen dentro de un círculo de radio $s$ desde el origen es:
\begin{equation}
    y = a \cdot \ln f(s) + c
    \label{eq:stouffer_integral}
\end{equation}
donde $f(s)$ es el número acumulado de oportunidades dentro de ese radio y $c$ es una constante de integración. Stouffer verificó empíricamente esta formulación con datos de movilidad residencial en Cleveland, Ohio (1933-1935). Los resultados muestran una buena correspondencia entre los flujos observados y los predichos por el modelo.

Liu y Yan \cite{liu2020intervening} formalizaron y extendieron esta familia de modelos de manera rigurosa. Para describir su funcionamiento, es necesario precisar tres conceptos fundamentales:

\begin{itemize}
    \item \textbf{Oportunidades de un lugar} ($m_l$): cantidad de oportunidades que ofrece el lugar $l$ para satisfacer el propósito del desplazamiento. En la práctica se asume que $m_l$ es proporcional a la población del lugar.
    \item \textbf{Oportunidades intervinientes} ($s_{ij}$): suma de oportunidades de todos los lugares cuya distancia al origen $i$ es menor que $d_{ij}$, excluyendo el propio origen $i$ y el destino $j$. Si se incluyen origen y destino se utiliza la notación $S_{ij} = m_i + s_{ij} + m_j$.
    \item \textbf{Valor de beneficio de una oportunidad} ($z$): utilidad o recompensa que una oportunidad reporta al viajero. Se modela como variable aleatoria continua con distribución $p(z)$.
\end{itemize}

El modelo de Stouffer establece que cada individuo que parte del origen $i$ ordena todos los destinos potenciales por distancia creciente y se detiene en cada lugar con una probabilidad fija proporcional al parámetro $\alpha$. Bajo este mecanismo, la probabilidad de que un individuo en $i$ elija el destino $j$ adopta la forma \cite{stouffer1940intervening, liu2020intervening}:
\begin{equation}
    P_{ij} = \frac{e^{-\alpha(S_{ij} - m_j)} - e^{-\alpha S_{ij}}}{1 - e^{-\alpha M}}
    \label{eq:io_formal}
\end{equation}
donde $m_j$ es la cantidad de oportunidades del destino $j$, $S_{ij}$ es el total de oportunidades en el círculo de radio $d_{ij}$ centrado en $i$ (incluyendo origen y destino), $\alpha$ es el parámetro de absorción a calibrar con datos históricos, y $M = \sum_k m_k$ es el total de oportunidades del sistema. La ecuación \eqref{eq:io_formal} muestra que la probabilidad de elegir $j$ crece con las oportunidades del destino $m_j$ y decrece con el número de oportunidades intervinientes $s_{ij}$: cuanto mayor es la competencia en el camino, menor la probabilidad de alcanzar el destino.

La principal limitación del modelo original de Stouffer es que contiene el parámetro $\alpha$, que debe calibrarse mediante datos históricos de flujo observado. Esta dependencia paramétrica restringe su aplicabilidad cuando no se dispone de matrices OD previas, y motivó el desarrollo de modelos libres de parámetros basados en el mismo concepto de oportunidades intervinientes.

\subsubsection{Modelo de radiación}

Simini et al. \cite{simini2012universal} propusieron en 2012 el modelo de radiación como una alternativa libre de parámetros al modelo gravitatorio. El modelo establece una analogía con los procesos de radiación y absorción en física de partículas. Su mecanismo se basa en un proceso estocástico de toma de decisiones locales: cada individuo en un origen busca la mejor oportunidad disponible dentro de radios geográficos crecientes y se desplaza hacia el destino donde la encuentra.

El mecanismo es el siguiente. Se asigna a cada oportunidad un valor de beneficio $z$ extraído de una distribución de probabilidad $p(z)$. La probabilidad de que un individuo en la zona $i$ elija la zona $j$ como destino es la probabilidad de que el mejor beneficio dentro de $j$ supere al mejor beneficio disponible dentro del círculo de radio $d_{ij}$ centrado en $i$ (excluyendo $i$ y $j$). Esta probabilidad es:
\begin{equation}
    p(n_j \mid m_i, s_{ij}) = \frac{m_i \, n_j}{(m_i + s_{ij})(m_i + n_j + s_{ij})}
    \label{eq:radiacion}
\end{equation}
donde $m_i$ es la población de la zona de origen $i$, $n_j$ es la población de la zona de destino $j$, y $s_{ij}$ es la población total dentro de un círculo de radio $d_{ij}$ centrado en $i$, excluyendo las poblaciones de $i$ y de $j$. El flujo esperado entre $i$ y $j$ es entonces:
\begin{equation}
    \langle T_{ij} \rangle = T_i \cdot \frac{m_i \, n_j}{(m_i + s_{ij})(m_i + n_j + s_{ij})}
    \label{eq:radiacion_flujo}
\end{equation}
donde $T_i = \sum_j T_{ij}$ es el flujo total saliente de la zona $i$. Una propiedad destacable del modelo de radiación es que satisface automáticamente la restricción de conservación de flujos: $\sum_j \langle T_{ij} \rangle = T_i$, sin necesidad de calibración de parámetros adicionales.

Simini et al. \cite{simini2012universal} validaron el modelo sobre datos de \textit{commuting} del censo de Estados Unidos (2000). Los resultados muestran que el modelo de radiación supera al modelo gravitatorio en la reproducción de flujos observados, especialmente en viajes de corta distancia, que el modelo gravitatorio tiende a sobreestimar.

\paragraph{Variantes del modelo de radiación.} La solidez del modelo original motivó diversas extensiones para adaptar su mecanismo a diferentes contextos. Ren et al. \cite{ren2014predicting} propusieron una versión basada en tiempo de viaje en lugar de distancia euclidiana, incorporando así la resistencia real de la red de transporte. Kang et al. \cite{kang2015generalized} formularon un modelo de radiación generalizado que introduce parámetros de escala y restricciones explícitas de producción y atracción, mejorando la precisión en contextos con heterogeneidad espacial marcada. Yang et al. \cite{yang2014limits} extendieron el modelo para incorporar la distribución no homogénea de oportunidades en el espacio, relevante en regiones con alta segregación funcional del suelo.

\paragraph{Limitación del modelo de radiación a escala intraurbana.} A pesar de su éxito en la predicción de \textit{commuting} interurbano, el modelo de radiación presenta una limitación estructural cuando se aplica a la escala intraurbana \cite{yan2014population, liu2020intervening}. Su supuesto central, que cada individuo elige el destino más cercano cuyo beneficio supera al del origen, refleja adecuadamente el comportamiento en la búsqueda de empleo a larga distancia (donde la proximidad es costosa), pero no captura la lógica real de los desplazamientos dentro de una ciudad. En entornos urbanos, las personas consideran simultáneamente el conjunto de destinos disponibles en toda la ciudad, sin restringirse a los más cercanos. Estudios empíricos sobre datos de transporte público en diversas ciudades confirman que el modelo de radiación subestima sistemáticamente los flujos intraurbanos de largo alcance \cite{yan2014population}. Adicionalmente, el modelo utiliza exclusivamente la población como \textit{proxy}\footnote{Una variable \textit{proxy} es una medida indirecta que se emplea en sustitución de una variable inobservable, asumiendo que existe una fuerte correlación estadística entre ambas.} de oportunidades, sin incorporar características del tejido urbano como el uso de suelo o la distribución de actividades económicas. Estas limitaciones fueron el punto de partida para el desarrollo del modelo de oportunidades ponderadas por población, descrito en la Sección \ref{sec:modelos_evolucion}, y justifican el uso de modelos basados en aprendizaje profundo capaces de representar el tejido urbano más allá de la simple distribución poblacional, como \textit{Deep Gravity}.

% ----------------------------------------------------------------
\subsection{Aprendizaje profundo aplicado a la movilidad}
% ----------------------------------------------------------------

El aprendizaje profundo (\textit{deep learning}) es un subcampo del aprendizaje automático que utiliza redes neuronales con múltiples capas de transformaciones no lineales para aprender representaciones jerárquicas de los datos \cite{liu2024interdisciplinary}. Su aplicación a la predicción de flujos de movilidad ha permitido incorporar un conjunto amplio de variables geoespaciales y capturar relaciones no lineales que los modelos clásicos no pueden representar.

Una red neuronal \textit{feedforward} (de propagación hacia adelante) con $L$ capas ocultas transforma un vector de entrada $\mathbf{x} \in \mathbb{R}^d$ mediante una secuencia de transformaciones afines seguidas de funciones de activación no lineales:
\begin{equation}
    \mathbf{z}^{(0)} = \sigma\!\left(W^{(0)}\,\mathbf{x}\right), \qquad
    \mathbf{z}^{(\ell)} = \sigma\!\left(W^{(\ell)}\,\mathbf{z}^{(\ell-1)}\right), \quad \ell = 1, \dots, L-1
    \label{eq:feedforward}
\end{equation}
donde $W^{(\ell)}$ es la matriz de pesos de la capa $\ell$ y $\sigma(\cdot)$ es la función de activación, comúnmente ReLU ($\max(0, x)$) o su variante LeakyReLU ($\max(\alpha x, x)$ con $\alpha$ pequeño), que mitiga el problema del gradiente desvanecido en redes profundas.

En el contexto de la generación de flujos, el aprendizaje profundo ha sido aplicado de múltiples formas: como extensión no lineal del modelo gravitatorio \cite{simini2021deep}, mediante redes neuronales de grafos que modelan la estructura espacial de las zonas, y mediante modelos de series de tiempo que capturan la evolución temporal de los flujos \cite{luca2024tsmob}. Esta memoria se enfoca en el primero de estos enfoques, implementado a través del modelo \textit{Deep Gravity}.

% ----------------------------------------------------------------
\subsection{Deep Gravity}
% ----------------------------------------------------------------

\textit{Deep Gravity} es un modelo de generación de flujos de movilidad propuesto por Simini et al. \cite{simini2021deep} en 2021. Su contribución central es demostrar que el modelo gravitatorio clásico restringido en origen (ecuación \ref{eq:gravedad_restringida}) es formalmente equivalente a una red neuronal lineal de una sola capa con salida \textit{softmax}, y que reemplazar esa capa lineal por una red neuronal profunda permite capturar relaciones no lineales entre las características del territorio y los flujos observados.

\subsubsection{Formulación matemática}

El modelo gravitatorio restringido en origen puede reescribirse como un modelo lineal generalizado con distribución multinomial. Para un origen $i$, la probabilidad de que un viaje se dirija al destino $j$ es:
\begin{equation}
    p_{ij} = \frac{\exp\!\left(\boldsymbol{\beta}^{\top}\,\mathbf{x}(l_i, l_j)\right)}{\displaystyle\sum_{k} \exp\!\left(\boldsymbol{\beta}^{\top}\,\mathbf{x}(l_i, l_k)\right)}
    \label{eq:dg_glm}
\end{equation}
donde $\mathbf{x}(l_i, l_j) = \bigl[\ln P_j,\, d_{ij}\bigr]$ es el vector de características del par $(i, j)$, y $\boldsymbol{\beta} = [\beta_1, \beta_2]$ son los parámetros del modelo.

Con distancia sin transformar, el numerador es proporcional a $P_j^{\beta_1}e^{\beta_2 d_{ij}}$; la disuasión exponencial corresponde a $\beta_2<0$. Una ley de potencia se obtiene incorporando $\ln d_{ij}$ en lugar de $d_{ij}$, para distancias positivas. Ambas son formulaciones gravitatorias, con funciones de disuasión distintas.

La formulación publicada utiliza la siguiente entropía cruzada entre distribuciones de destino normalizadas por origen, para orígenes con $O_i>0$:
\begin{equation}
    H = -\sum_i \sum_j \frac{T_{ij}}{O_i} \ln p_{ij}
    \label{eq:dg_loss}
\end{equation}
Esta expresión asigna el mismo peso a cada origen. La pérdida efectiva de Santiago pondera las contribuciones por la masa observada, $-\sum_i\sum_j T_{ij}\ln p_{ij}$, como se especifica en la ecuación~\eqref{eq:perdida_multinomial_santiago}. Las dos expresiones difieren en la ponderación entre orígenes y no se presentan como objetivos de entrenamiento idénticos.
\textit{Deep Gravity} reemplaza el producto interno lineal $\boldsymbol{\beta}^{\top}\,\mathbf{x}(l_i, l_j)$ por una red neuronal profunda $f_\theta(\mathbf{x}(l_i, l_j))$ que produce un \textit{score} escalar. Las probabilidades de destino se obtienen mediante \textit{softmax} sobre los \textit{scores} de todos los destinos candidatos desde el origen $i$:
\begin{equation}
    p_{ij} = \frac{\exp\!\left(f_\theta\!\left(\mathbf{x}(l_i, l_j)\right)\right)}{\displaystyle\sum_{k} \exp\!\left(f_\theta\!\left(\mathbf{x}(l_i, l_k)\right)\right)}
    \label{eq:dg_softmax}
\end{equation}
El flujo estimado desde $i$ hacia $j$ es entonces $\hat{T}_{ij} = O_i \cdot p_{ij}$. El total $O_i$ se considera conocido; la tarea predictiva consiste en estimar la distribución de ese flujo entre destinos.

\subsubsection{Arquitectura de la red neuronal}

En la arquitectura descrita por Simini et al. \cite{simini2021deep}, la función $f_\theta$ se implementa como una red neuronal \textit{feedforward} de 15 capas ocultas con activación LeakyReLU. Las primeras 6 capas tienen dimensión 256 y las 9 restantes dimensión 128. La capa de salida produce un escalar:
\begin{align}
    \mathbf{z}^{(0)} &= \text{LeakyReLU}\!\left(W^{(0)}\,\mathbf{x}(l_i, l_j)\right) \nonumber \\
    \mathbf{z}^{(\ell)} &= \text{LeakyReLU}\!\left(W^{(\ell)}\,\mathbf{z}^{(\ell-1)}\right), \quad \ell = 1, \dots, 14 \nonumber \\
    f_\theta\!\left(\mathbf{x}(l_i, l_j)\right) &= \mathbf{w}^{\top}\,\mathbf{z}^{(14)}
    \label{eq:dg_red}
\end{align}
Esta distribución de anchos corresponde a la arquitectura descrita en la publicación. La implementación utilizada en Santiago tiene cinco capas ocultas de 256 unidades y diez de 128; su configuración efectiva se documenta en la Sección~\ref{sec:configuracion_entrenamiento}.
El protocolo publicado emplea RMSprop, con momento 0,9, tasa de aprendizaje $5\times10^{-6}$, lotes de 64 orígenes y hasta 512 destinos muestreados por origen. El muestreo permite calcular la normalización sobre un subconjunto de candidatos durante el entrenamiento; no equivale a evaluar todos los destinos. La configuración efectivamente usada en el piloto de Santiago se documenta en la Sección~\ref{sec:configuracion_entrenamiento}: emplea lotes de ocho orígenes y el soporte completo de destinos. Los parámetros de los experimentos finales se fijarán en su diseño específico.

\subsubsection{Vector de características geoespaciales}

En la formulación de referencia \cite{simini2021deep}, el vector de entrada reúne atributos del origen y del destino, población y distancia. Su estructura puede representarse mediante el siguiente esquema de 39 componentes:
\begin{equation}
    \mathbf{x}(l_i, l_j) = \bigl[\mathbf{f}(l_i),\; \mathbf{f}(l_j),\; \ln P_i,\; \ln P_j,\; d_{ij}\bigr]
    \label{eq:dg_features}
\end{equation}
donde $\mathbf{f}(l)\in\mathbb{R}^{18}$ reúne atributos geoespaciales de la ubicación $l$, $P_i$ y $P_j$ representan la población asociada a ambos extremos y $d_{ij}$ es la distancia entre sus representantes espaciales. La procedencia de las variables y sus transformaciones deben identificarse en cada aplicación.

En Santiago se utiliza un conjunto de 19 atributos por ubicación y 39 componentes por par. Incluye densidad de población censal, coberturas de suelo, densidades viales y atributos de equipamientos y servicios. Su construcción, orden y transformación se definen en la Sección~\ref{sec:features_osm}. La igualdad de dimensión con el vector de referencia no implica equivalencia de los atributos. La población de Santiago procede del Censo 2024 y se expresa como densidad después de la interpolación areal; no se introduce directamente el logaritmo de la población total ni el flujo saliente observado como variable de población.

La agregación espacial depende de la magnitud del atributo. En Santiago, la densidad de población se expresa en personas por km$^2$, la densidad de puntos de interés en puntos por km$^2$ y la densidad vial en km de vías por km$^2$. La cobertura corresponde al área de la unión de polígonos de una categoría dividida por el área de la unidad, por lo que es una proporción adimensional entre cero y uno. Estas operaciones se distinguen de las transformaciones y la estandarización previas al entrenamiento, descritas en la Sección~\ref{sec:features_osm}.

\subsubsection{Resultados y desempeño}

Simini et al. \cite{simini2021deep} evaluaron \textit{Deep Gravity} sobre flujos censales de \textit{commuting} en Inglaterra e Italia y flujos derivados de telefonía móvil en el estado de Nueva York, utilizando la métrica \textit{Common Part of Commuters} (CPC), definida como:
\begin{equation}
    \text{CPC}(\hat{\mathbf{T}}, \mathbf{T}) = \frac{2\,\displaystyle\sum_{i,j} \min\!\left(\hat{T}_{ij},\, T_{ij}\right)}{\displaystyle\sum_{i,j} \hat{T}_{ij} + \displaystyle\sum_{i,j} T_{ij}}
    \label{eq:cpc}
\end{equation}
Cuando la masa del denominador es positiva, la métrica CPC toma valores en $[0,1]$, donde 1 indica coincidencia exacta entre flujos generados y reales. Es equivalente al índice de Sørensen-Dice aplicado a flujos. En los experimentos publicados, \textit{Deep Gravity} supera al modelo gravitatorio clásico en todos los contextos evaluados. Las mejoras en CPC global alcanzan hasta un 190\% en Inglaterra y un 1.066\% en el estado de Nueva York. El estudio también evaluó la transferencia geográfica. Entrenado en algunas ciudades del Reino Unido, mantuvo un desempeño similar al evaluarse en ciudades no vistas durante el entrenamiento. En el caso de Nueva York, la publicación aproxima la población mediante flujo saliente observado; esa decisión no se traslada a Santiago. La auditoría técnica del código no acredita por sí sola una reproducción integral de esos resultados.

\subsubsection{Síntesis comparativa de modelos}

La Tabla~\ref{tab:comparacion_modelos} resume las formulaciones examinadas en la literatura. Es una comparación bibliográfica, no una tabla de resultados experimentales de esta memoria. La calibración, los atributos y el comportamiento intraurbano dependen de la variante; la evidencia latinoamericana se limita a las fuentes revisadas.

\begin{table}[p]
\centering
\caption{Comparación bibliográfica de modelos de flujos. Fuente: elaboración propia a partir de las referencias indicadas.}
\label{tab:comparacion_modelos}
\footnotesize
\renewcommand{\arraystretch}{1.15}
\begin{tabularx}{\textwidth}{@{}p{2.4cm}XXXXX@{}}
\toprule
\textbf{Modelo} &
\textbf{Atributos considerados} &
\textbf{Calibración} &
\textbf{Escala intraurbana} &
\textbf{Interpretación de la predicción} &
\textbf{Evidencia latinoamericana} \\
\midrule
Gravitatorio clásico \cite{zipf1946p,liu2024interdisciplinary}
    & Población y distancia
    & Exponentes y disuasión
    & Evaluación según contexto
    & Formulación explícita
    & No documentada aquí \\
Oportunidades intervinientes \cite{stouffer1940intervening,liu2020intervening}
    & Oportunidades y su distribución
    & Parámetro de absorción
    & Depende de la formulación
    & Formulación explícita
    & No documentada aquí \\
Radiación original \cite{simini2012universal,yan2014population}
    & Población como aproximación de oportunidades
    & Sin parámetros de forma
    & Limitaciones documentadas
    & Formulación explícita
    & No documentada aquí \\
Variantes de oportunidades intervinientes \cite{yan2014population,liu2019opportunity}
    & Oportunidades, habitualmente aproximadas por población
    & Depende de la variante
    & Evaluadas en los contextos citados
    & Formulación explícita
    & No documentada aquí \\
\textit{Deep Gravity} \cite{simini2021deep}
    & Población, atributos territoriales y distancia
    & Pesos de la red
    & Evaluación según contexto
    & Atribución posterior, como SHAP
    & No documentada aquí \\
\bottomrule
\end{tabularx}

\smallskip
\begin{minipage}{\textwidth}
\footnotesize
La última columna significa que no se documenta una validación de esa formulación en las fuentes utilizadas para esta comparación; no demuestra la inexistencia de otros estudios. Una formulación explícita facilita inspeccionar sus supuestos, pero no garantiza una explicación causal. SHAP incluye aproximaciones agnósticas al modelo y no es exclusivo de las redes profundas \cite{lundberg2017shap}. El artículo original de \textit{Deep Gravity} ya incorpora un análisis SHAP; su aplicación a Santiago permanece pendiente.
\end{minipage}
\end{table}

La combinación de atributos territoriales explícitos, una función predictiva no lineal y la posibilidad de aplicar métodos de atribución motiva la evaluación de \textit{Deep Gravity} en esta memoria. La comparación bibliográfica no establece por sí sola su superioridad en Santiago ni demuestra que estas capacidades sean exclusivas de esa arquitectura.

% ----------------------------------------------------------------
\subsection{Explicabilidad en modelos de inteligencia artificial (XAI)}
% ----------------------------------------------------------------

La explicabilidad de la inteligencia artificial (\textit{eXplainable Artificial Intelligence}, XAI) se refiere al conjunto de técnicas que permiten comprender e interpretar las predicciones de modelos de aprendizaje automático complejos \cite{lundberg2017shap}. En el contexto de la planificación del transporte, la explicabilidad es relevante no solo desde el punto de vista académico, para estudiar cómo las variables de entrada contribuyen a las predicciones, sino también desde el punto de vista institucional, ya que permite a los planificadores auditar las predicciones del modelo.

\subsubsection{Valores SHAP}

SHAP (\textit{SHapley Additive exPlanations}), propuesto por Lundberg y Lee \cite{lundberg2017shap}, es un método de explicabilidad que fundamenta la importancia de cada variable en la teoría de juegos cooperativos de Shapley. La contribución de la característica $i$ a la predicción del modelo para una instancia $\mathbf{x}$ se define como:
\begin{equation}
    \phi_i(f, \mathbf{x}) = \sum_{S \subseteq \mathcal{F} \setminus \{i\}} \frac{|S|!\,(|\mathcal{F}|-|S|-1)!}{|\mathcal{F}|!}\,\bigl[f(S \cup \{i\}) - f(S)\bigr]
    \label{eq:shap}
\end{equation}
donde $\mathcal{F}$ es el conjunto de todas las características, $S$ es una coalición de características que no incluye a $i$, y $f(S)$ es la predicción del modelo usando solo las características en $S$. El valor $\phi_i$ representa la contribución marginal promedio de la característica $i$ sobre todas las posibles coaliciones de características.

Los valores SHAP satisfacen tres propiedades deseables: \textit{eficiencia} (la suma de todos los valores SHAP es igual a la diferencia entre la predicción y el valor base), \textit{simetría} (características con igual contribución marginal reciben igual valor SHAP) y \textit{dummy} (características sin influencia reciben valor cero). En conjunto, estas propiedades garantizan una atribución justa y consistente de importancias.

Los valores SHAP permiten estudiar la contribución de los atributos a una predicción con respecto a una referencia. Su aplicación al modelo final de Santiago se contempla como análisis posterior a la selección experimental. La configuración y la estabilidad de ese análisis se documentarán en el capítulo de resultados; las atribuciones se interpretarán como propiedades de las predicciones, sin inferir efectos causales de intervenciones urbanas.

\subsubsection{Gradientes integrados}

Los \textit{Integrated Gradients} \cite{sundararajan2017axiomatic} son un método de atribución basado en el gradiente de la función del modelo con respecto a las variables de entrada. A diferencia de SHAP, que se fundamenta en la teoría de juegos cooperativos, los \textit{Integrated Gradients} derivan las atribuciones directamente de la estructura diferencial de la red neuronal. Para una predicción $f(\mathbf{x})$ y una entrada de referencia $\mathbf{x}'$, la atribución de la característica $i$ se define como:
\begin{equation}
    \text{IG}_i(\mathbf{x}) = (x_i - x_i') \cdot \int_0^1 \frac{\partial f\!\left(\mathbf{x}' + \alpha(\mathbf{x} - \mathbf{x}')\right)}{\partial x_i}\,\mathrm{d}\alpha
    \label{eq:ig}
\end{equation}
El método satisface dos propiedades axiomáticas que lo distinguen de otras técnicas de atribución: \textit{sensibilidad} (si la entrada y la referencia difieren en una sola característica y producen predicciones distintas, esa característica recibe atribución no nula) e \textit{invariancia a la implementación} (la atribución depende únicamente de la función matemática que computa la red, no de su arquitectura interna). Estas propiedades lo hacen complementario a SHAP, que satisface eficiencia, simetría y dummy.

La elección de la entrada de referencia $\mathbf{x}'$ es un aspecto crítico del método. El vector nulo ($\mathbf{x}' = \mathbf{0}$) es una convención habitual en clasificación de imágenes, pero su significado depende de las variables y de su preprocesamiento. Una densidad cartografiada puede ser cero; después de estandarizar, el valor cero representa la media de entrenamiento. La referencia debe definirse en el espacio de entrada del modelo y su elección requiere un análisis de sensibilidad. El protocolo explicativo determinará la referencia que se utilizará en Santiago.

Los gradientes integrados pueden considerarse una técnica complementaria para atribuir una salida diferenciable del modelo a sus entradas. Su eventual aplicación en Santiago dependerá del protocolo de explicabilidad del modelo final.

% ----------------------------------------------------------------
\subsection{OpenStreetMap como fuente de datos geoespaciales}
% ----------------------------------------------------------------

OpenStreetMap (OSM) es una base de datos geoespacial colaborativa y de libre acceso, construida y mantenida por una comunidad global de voluntarios desde su creación en 2004 \cite{haklay2008openstreetmap}. Su modelo de datos se articula en tres tipos de objetos geométricos: \textit{nodos} (puntos), \textit{vías} (secuencias de nodos que representan líneas o polígonos) y \textit{relaciones} (agrupaciones de nodos y vías con un significado semántico compartido). Cada objeto puede llevar asociados pares clave-valor denominados \textit{etiquetas} (\textit{tags}), que codifican su naturaleza (por ejemplo, \texttt{amenity=hospital} o \texttt{landuse=residential}).

Para modelar flujos de movilidad, OpenStreetMap aporta atributos de uso del suelo, vialidad, equipamientos y servicios. En esta memoria, el término \textit{amenities} se refiere a equipamientos y servicios urbanos; los atributos de OSM abarcan también información que no pertenece a ese conjunto. Los puntos de interés son una representación de algunas de estas entidades, mientras que otras se describen mediante áreas o líneas. La selección y construcción de los atributos empleados en Santiago se describe en la Sección~\ref{sec:features_osm}.

Una limitación reconocida de OSM en el contexto latinoamericano es la heterogeneidad espacial de su cobertura. Estudios previos \cite{BarringtonLeigh2017} han documentado que la densidad y completitud de los datos OSM varía significativamente según el nivel socioeconómico y la accesibilidad de las distintas áreas de una ciudad. En Santiago, esto se manifiesta en una mayor completitud en comunas de ingresos medios y altos, como Providencia, Las Condes y Vitacura, y menor cobertura en comunas periféricas de menores ingresos. Esta heterogeneidad es un factor metodológico relevante para esta memoria. Puede sesgar la estimación de las características geoespaciales de entrada al modelo en las zonas con menor presencia de datos voluntarios. Este aspecto se aborda en el análisis de resultados.

% ----------------------------------------------------------------
\subsection{Contexto de la movilidad en Santiago de Chile}
% ----------------------------------------------------------------

El sistema integrado de transporte público de Santiago reúne buses, Metro y el servicio ferroviario Nos-Estación Central \cite{dtpm_informe_2024}. Su cobertura operativa y las distintas delimitaciones estadísticas de la ciudad no deben confundirse con el área de estudio, definida aquí como un núcleo urbano referencial de 34 comunas. La delimitación utilizada se describe en la Sección~\ref{sec:delimitacion}.

Las limitaciones de las fuentes de datos de movilidad disponibles para Santiago (la EOD de SECTRA y los registros de validación de la tarjeta bip!) se detallan en la Sección~\ref{sec:fuentes_movilidad}.

Esta memoria utiliza la entrega DTPM disponible entre el 1 y el 29 de noviembre de 2024. La unidad de análisis se reconstruye al unir viajes y etapas: las zonas y el factor de expansión principal provienen de viajes, mientras que las coordenadas UTM de subida y bajada se obtienen de etapas. La fuente permite construir matrices OD diarias y por período horario; los archivos del día 30 se excluyen por errores, sin imputación. Esta entrega constituye la principal fuente de flujos observados para el entrenamiento y la evaluación de \textit{Deep Gravity} en Santiago.
